\documentclass[11pt]{amsart}
\usepackage[T1]{fontenc}
\usepackage{amsmath,amssymb,amsthm,mathtools}
\usepackage{xcolor}
\usepackage[colorlinks=true,linkcolor=blue!55!black,citecolor=blue!55!black,urlcolor=blue!55!black]{hyperref}

\newtheorem{theorem}{Theorem}[section]
\newtheorem{lemma}[theorem]{Lemma}
\newtheorem{proposition}[theorem]{Proposition}
\newtheorem{corollary}[theorem]{Corollary}
\theoremstyle{definition}
\newtheorem{definition}[theorem]{Definition}
\newtheorem{remark}[theorem]{Remark}

\title[Shifted-Golomb Prefix Identity]{A Shifted-Golomb Identity for the Full-Memory Prefix of Run-Constructed Almost Golomb Sequences}
\author{Zijian Zeng}
\address{Institute of Computer Science and Digital Innovation, UCSI University, Kuala Lumpur 56000, Malaysia}
\email{zijianzeng@foxmail.com}
\email{1002266693@ucsiuniversity.edu.my}
\date{2026-08-29}
\hypersetup{pdftitle={A Shifted-Golomb Identity for the Full-Memory Prefix of Run-Constructed Almost Golomb Sequences},pdfauthor={Zijian Zeng}}

\begin{document}

\begin{abstract}
For the earliest, run-constructed almost Golomb sequence $a_r$ of order
$r\geq3$, we prove the stronger prefix identity
$a_r(n)=G(n-1)$ for every $3\leq n\leq r$, where $G$ is the classical
Golomb sequence.  In particular, this proves the Prefix Conjecture
$a_r(r)=G(r-1)$.  The proof identifies the exact conjugacy between the
cumulative run endpoints of the two sequences.  We also make explicit a
necessary interpretive repair: the pointwise minimality sentence in the
original definition is inconsistent at $n=2$ if read literally as a
condition on the completed sequence, so our theorem uses the run
construction stated in the source's existence theorem.  Exact computations
with two independent implementations are reported only as diagnostics.
\end{abstract}

\maketitle

\section{Introduction}

Golomb's self-describing sequence $G$ is the unique nondecreasing sequence
of positive integers in which the value $m$ occurs exactly $G(m)$ times
\cite{Golomb1966}.  Cloitre introduced finite-memory analogues by replacing
the full cumulative sum with a sliding window of length $r$ and selecting an
earliest monotone solution \cite{Cloitre2026Almost}.  If $a_r(k)=0$ for
$k\leq0$, the implicit equation is
\begin{equation}\label{eq:implicit}
 a_r\!\left(\sum_{j=0}^{r-1}a_r(n-j)\right)=n.
\end{equation}
The source conjectures that
\begin{equation}\label{eq:source-conjecture}
 a_r(r)=G(r-1)\qquad(r\geq3)
\end{equation}
and reports a computational verification for $r\leq200$
\cite[Conjecture~9.3]{Cloitre2026Almost}.

We prove more: every term from index $3$ through index $r$ has the same
shifted-Golomb form.  The proof is a run-endpoint conjugacy.  The only
subtlety is definitional.  As explained in Section~\ref{sec:definition}, the
pointwise formula used to express minimality in the source cannot be read
literally at $n=2$.  We therefore take the source's run construction
\cite[Theorem~2.2]{Cloitre2026Almost} as the precise normalization.  This
also excludes nongreedy solutions of the bare implicit equation, whose
existence is studied separately in \cite{Cloitre2026Beatty}.

\section{Definition repair and the run normalization}\label{sec:definition}

The source supplements \eqref{eq:implicit} by saying that no smaller
candidate $m$ in a specified interval satisfies the same equation.  If this
is interpreted as a condition on the completed sequence, then
\eqref{eq:implicit} forces $a_r(1)=1$ and $a_r(2)=2$, while the smaller
candidate $m=1$ at $n=2$ gives

\[
 a_r(m+a_r(1))=a_r(2)=2.
\]

Thus the literal pointwise reading is inconsistent.  The intended object is
nevertheless clear from the run construction in the source.  We record it
non-circularly.

\begin{definition}\label{def:run}
Fix $r\geq3$, put $a_r(k)=0$ for $k\leq0$, and initialize
\[
 a_r(1)=1,\qquad a_r(2)=a_r(3)=a_r(4)=2,\qquad P_3=5.
\]
At stage $m\geq3$, set
\begin{align*}
 L_m&=a_r(m+1)-a_r(m+1-r),\\
 a_r(q)&=m\quad(P_m\leq q<P_m+L_m),\\
 P_{m+1}&=P_m+L_m.
\end{align*}
\end{definition}

\begin{lemma}\label{lem:well-founded}
Definition~\ref{def:run} is well founded and defines an infinite
nondecreasing sequence.  At stage $m$ one has $P_m\geq m+2$ and
$1\leq L_m\leq r$.
\end{lemma}

\begin{proof}
At stage $3$, the defined prefix is $(1,2,2,2)$, so the assertion follows
from $P_3=5$ and $r\geq3$.  Suppose the assertion holds at stage $m$.
Then $m+1<P_m$, so $a_r(m+1)$ is already defined.  The extended sequence,
with value zero at nonpositive indices, has increments in $\{0,1\}$.
Consequently $0\leq L_m\leq r$.

If $L_m=0$, monotonicity makes $a_r$ constant on the $r+1$ consecutive
indices from $m+1-r$ through $m+1$.  The left endpoint cannot be
nonpositive, because its value would then be zero while $a_r(m+1)>0$.
Hence these $r+1$ positions lie in one previously completed run, contrary
to the inductive bound that every such run has length at most $r$.
Thus $1\leq L_m\leq r$.  Appending $L_m$ copies of $m$ preserves unit
increments, completes the next run, and gives
$P_{m+1}=P_m+L_m\geq m+3$.  Induction proves the claim, and the positive
run lengths imply $P_m\to\infty$.
\end{proof}

For later use put
\[
 S_m^{(r)}=\sum_{j=0}^{r-1}a_r(m-j).
\]

\begin{lemma}\label{lem:endpoints}
For every $m\geq3$, $P_m=S_m^{(r)}$.  Hence the run of $m$ in $a_r$ is
exactly
\[
 [S_m^{(r)},S_{m+1}^{(r)}-1].
\]
\end{lemma}

\begin{proof}
The base identity is $S_3^{(r)}=2+2+1=5=P_3$.  Moreover,
\[
 S_{m+1}^{(r)}-S_m^{(r)}
 =a_r(m+1)-a_r(m+1-r)=L_m=P_{m+1}-P_m.
\]
The result follows by induction and Definition~\ref{def:run}.
\end{proof}

\section{Shifted Golomb endpoints}

Set
\[
 H(0)=0,\qquad H(m)=\sum_{j=1}^{m}G(j)\quad(m\geq1).
\]
The defining multiplicity property says exactly that the run of $m$ in $G$
is $[H(m-1)+1,H(m)]$.  Define the shifted sequence
\[
 b(1)=1,\qquad b(2)=2,\qquad b(n)=G(n-1)\quad(n\geq3)
\]
and let $T_m=\sum_{j=1}^{m}b(j)$.

\begin{lemma}\label{lem:shift}
For every $m\geq3$,
\[
 T_m=H(m-1)+2,
\]
and the run of $m$ in $b$ is $[T_m,T_{m+1}-1]$.
\end{lemma}

\begin{proof}
The sum identity is
\[
 T_m=1+2+\sum_{j=2}^{m-1}G(j)=H(m-1)+2.
\]
The condition $b(q)=m$ is equivalent to $G(q-1)=m$.  Shifting the Golomb
run by one therefore gives
\[
 [H(m-1)+2,H(m)+1]=[T_m,T_{m+1}-1].
\]
\end{proof}

\section{The strong prefix theorem}

\begin{theorem}\label{thm:main}
For every integer $r\geq3$ and every $3\leq n\leq r$,
\[
 a_r(n)=G(n-1).
\]
In particular, the Prefix Conjecture \eqref{eq:source-conjecture} holds.
\end{theorem}

\begin{proof}
We prove $a_r(n)=b(n)$ by strong induction.  The definitions agree through
position $4$: both prefixes are $(1,2,2,2)$.

Suppose $5\leq n\leq r$ and equality holds at all smaller positive indices.
Put $m=b(n)=G(n-1)$.  We have $m\geq3$ because $G(4)=3$.  Moreover,
$m\leq n-2$: the nondecreasing sequence $G$ cannot skip a positive integer,
and its value $2$ is repeated at positions $2$ and $3$, so
$G(k)\leq k-1$ for every $k\geq3$.  Thus $m+1<n$.

Since $m+1\leq n-1<r$, both relevant sliding windows contain every positive
index through their endpoints.  Every term involved also has index below
$n$, so the induction hypothesis gives
\begin{align*}
 S_m^{(r)}&=\sum_{j=1}^{m}a_r(j)
            =\sum_{j=1}^{m}b(j)=T_m,\\
 S_{m+1}^{(r)}&=T_{m+1}.
\end{align*}
Lemma~\ref{lem:shift} places $n$ in $[T_m,T_{m+1}-1]$, while
Lemma~\ref{lem:endpoints} identifies that same interval as the run of $m$ in
$a_r$.  Hence $a_r(n)=m=b(n)$, completing the induction.
\end{proof}

\section{Exact computational diagnostics}

The preceding theorem is entirely symbolic.  As an adversarial indexing
check, we also used two exact-integer implementations.  The first constructs
$a_r$ from the telescoped run lengths and constructs $G$ by self-describing
run expansion; it checks every pair $(r,n)$ with $3\leq r\leq2000$ and
$3\leq n\leq r$.  The second recomputes both sliding sums directly and
constructs $G$ from its scalar recurrence; it checks through $r=500$.
Together they evaluate 2,121,252 prefix pairs, with no counterexample.
A separate verifier recomputes both outputs from the serialized finite
record.  These finite calculations are diagnostics, not proof of
Theorem~\ref{thm:main}.

\section{Equality cases and limitations}

The theorem covers every order in the source conjecture.  The cases $r=3$
and $r=4$ are contained in the explicit base prefix; the first inductive case
is $n=5$, where $m=3$ and both required endpoint indices are already known.
There are no limiting or zero-denominator cases.

The result concerns only the earliest/run-constructed branch.  It is not a
classification of all monotone solutions of \eqref{eq:implicit}, and it does
not address the separate Domination Conjecture or the global maximal
multiplicity proposed in \cite{Cloitre2026Almost}.

\section*{Data, code, and computational disclosure}

The accompanying source archive contains the two exact Python
implementations, the serialized finite diagnostic, and an independent
verifier.  They use Python 3.10.12, exact integers, no randomness, and no
external package.  The author reviewed the proof and takes responsibility
for all claims.  No proof assistant was used.

\begingroup
\raggedright
\bibliographystyle{amsplain}
\bibliography{references}
\endgroup

\end{document}
